\chapter{Oprogramowanie}

\section{Język Rust w programowaniu mikrokontrolerów}

\section{Narzędzia i biblioteki}


\section{Obsługa elementów pomiarowych}
\label{section:devices_programming}
W celu ograniczenia wykorzystania zasobów pamięciowych mikrokontrolera znaczna część konfiguracji urządzeń pomiarowych realizowana jest na etapie kompilacji. Rozbudowany system typów \translation{type system} języka Rust umożliwia zwięzłe i czytelne wyrażenie konfiguracji z wykorzystaniem typów generycznych \translation{generic types}, cech \translation{traits} oraz stałych czasu kompilacji \translation{const generics}.

\subsection{Moduł AHT20 --- czujnik temperatury i wilgotności względnej}
Komunikacja z modułem AHT20 realizowana jest za pośrednictwem magistrali \isqrc poprzez wysyłanie dedykowanych komend sterujących. Moduł nie wymaga przesyłania danych konfiguracyjnych w trakcie działania. Zgodnie z notą katalogową producenta \cite{AHT20} urządzenie obsługuje następujący zestaw komend:
\begin{itemize}
    \item \textbf{Inicjalizacja} --- wymuszenie procedury kalibracji wewnętrznej układu;
    \item \textbf{Wyzwolenie pomiaru} --- zainicjowanie akwizycji danych oraz
          ich zapis do wewnętrznej pamięci urządzenia;
    \item \textbf{Odczyt statusu} --- weryfikacja stanu kalibracji układu oraz
          gotowości do przeprowadzenia kolejnego pomiaru.
\end{itemize}

Procedura inicjalizacji modułu, przedstawiona na listingu \ref*{lst:aht20_init}, rozpoczyna się od odczekania czasu rozruchu wynoszącego \SI{40}{\milli\second}, po którym następuje odczyt rejestru statusu. W przypadku braku flagi kalibracji wysyłana jest komenda inicjalizacji, a następnie wprowadzane jest opóźnienie \SI{10}{\milli\second} niezbędne do jej wykonania:

\begin{lstlisting}[language=Rust, caption={Inicjalizacja modułu AHT20}, label={lst:aht20_init}]
delay.delay_ms(40);

let status = self.read_status(i2c)?;

if status & STATUS_CALIBRATED_BIT == 0 {
    i2c.write(self.address, &COMMAND_INITIALIZE)?;
    delay.delay_ms(10);
}
\end{lstlisting}


\subsection{Moduł BMP280}


\subsection{Moduł VEML7700}


\subsection{Modułu z licznikiem Geigera-Müllera}

\section{Komunikacja radiowa}
\subsection{Serializacja danych}
\subsection{Zabezpieczenie danych}
[Nie zaimplementowane]
\subsection{Transmisja danych}